% gl_pdf.tex — Growth Lab design system for R Markdown PDF output
%
% Usage in a YAML header (point in_header at this file wherever the kit lives):
%   output:
%     pdf_document:
%       latex_engine: xelatex
%       includes:
%         in_header: .../skills/gl-ggplot/assets/gl_pdf.tex
%
% Requires: xelatex (or lualatex). Variable fonts work in both engines.
% Fonts are resolved BY NAME via fontconfig, so the GL fonts must be installed
% on the system first (run scripts/install.sh, then `fc-cache -f` on Linux).
% This keeps the header portable — no machine-specific paths.
%
% Fonts (nil-design):
%   Inter             — main / sans / body (variable file)
%   Source Serif 4    — headings + display (registered as \glserif; variable)
%   JetBrains Mono    — code blocks only (kept in assets/fonts/legacy/)
%
% The variable-font opsz axis (Source Serif 4) is not addressed here; faces
% default to the text-cut shapes. See docs/followups.md #3.

% ---- Fonts ------------------------------------------------------------------

\usepackage{fontspec}

% Resolved by family name via fontconfig (see header note). Inter and Source
% Serif 4 ship as variable faces whose italic cut is a same-family sibling, so
% fontspec discovers upright + italic automatically; bold falls back to the
% renderer (heading weights 500/600 are not separately addressable here).
\setmainfont{Inter Variable}

\setsansfont{Inter Variable}

\newfontfamily\glserif{Source Serif 4 Variable}

\setmonofont{JetBrains Mono}[ Scale = 0.92 ]

% ---- Colours ----------------------------------------------------------------
% Mirrors the `gl` palette in theme_gl.R.

\usepackage{xcolor}
\definecolor{glink}{HTML}{1A1714}     % ink   — strong / headings / titles
\definecolor{glink2}{HTML}{2C2823}    % ink-2 — body
\definecolor{glink3}{HTML}{6B645A}    % ink-3 — captions, eyebrows, chrome
\definecolor{glink4}{HTML}{9A9389}    % ink-4 — hairlines

\definecolor{glaccent}{HTML}{015C9C}  % accent blue
\definecolor{glamber}{HTML}{C77A20}   % amber accent / story-highlight
\definecolor{glmuted}{HTML}{7E8A99}   % muted gray — "everyone-else" in charts

\definecolor{glrule}{HTML}{DDDDDD}    % hairline rules
\definecolor{glpaper}{HTML}{FAF8F4}   % page paper
\definecolor{glcodebg}{HTML}{F4F1EA}  % code-block tint

% Compatibility aliases (older docs may reference these)
\colorlet{glblue}{glaccent}
\colorlet{glred}{glamber}
\colorlet{gldark}{glink}
\colorlet{glborder}{glrule}
\colorlet{glbg}{glcodebg}

% ---- Headings ---------------------------------------------------------------
% Display title (cover) handled via the Word/Pandoc `\title` mechanism, not
% by section commands. Section commands here render H1 / H2 / H3.
% Heading weights 500/600 not directly representable — accept system bold.

\makeatletter
\renewcommand\section{%
  \@startsection{section}{1}{\z@}%
    {-3.0ex \@plus -1ex \@minus -.2ex}%
    {1.5ex \@plus.2ex}%
    {\glserif\LARGE\mdseries\color{glink}}}
\renewcommand\subsection{%
  \@startsection{subsection}{2}{\z@}%
    {-2.0ex \@plus -1ex \@minus -.2ex}%
    {0.6ex \@plus.2ex}%
    {\glserif\large\mdseries\color{glink}}}
\renewcommand\subsubsection{%
  \@startsection{subsubsection}{3}{\z@}%
    {-1.5ex \@plus -1ex \@minus -.2ex}%
    {0.5ex \@plus.2ex}%
    {\glserif\normalsize\bfseries\color{glink}}}  % followups #1 — H3 improvised
\makeatother

% ---- Links ------------------------------------------------------------------

\usepackage{hyperref}
\hypersetup{
  colorlinks = true,
  linkcolor  = glaccent,
  citecolor  = glaccent,
  urlcolor   = glaccent
}

% ---- Body text and spacing --------------------------------------------------

\setlength{\parskip}{0.6em}
\setlength{\parindent}{0pt}
\linespread{1.33}   % ~1.6 perceived line-height with xelatex's extra metrics

% ---- Captions ---------------------------------------------------------------
% Chart sources: italic serif distinguishes provenance from the data narrative.

\usepackage{caption}
\captionsetup{
  font       = {small, color=glink2},
  labelfont  = {bf, color=glaccent},
  textfont   = {it},
  labelsep   = colon,
  margin     = 4pt
}

% ---- Code blocks ------------------------------------------------------------
% Pandoc renders highlighted code in the `Shaded` environment with
% `\colorbox{shadecolor}{...}`. Override `shadecolor` to a soft warm grey.

\definecolor{shadecolor}{HTML}{F4F1EA}

% ---- Tables -----------------------------------------------------------------

\usepackage{booktabs}
\renewcommand{\arraystretch}{1.15}
